\section{Conditional loop}%
\label{sec:flow:while_loop}

Conditional loops continue iterating as long as a specified condition remains
true. They are initialized with the keyword \token{while} and follow the syntax
\token{while C V}, where \token{C} represents the condition that must be
satisfied, and \token{V} is the value evaluated while the condition remains
true. Similar to conditional branching, the condition \token{C} must be of type
\token{bool}. The following example showcases a while loop that iterates five
times before exiting. Its control flow is illustrated in
Figure~\ref{fig:flow:simple_while_loop}

\begin{lstlisting}[style=coloredverbatim, caption=Simple while loop, label=lst:flow:simple_while_loop]
let mut count = 0;

while count < 5 {
  count += 1;
}

assert(count == 5);
\end{lstlisting}

\begin{figure}[h]
  \centering
  \begin{adjustbox}{max width=\linewidth}
    \begin{tikzpicture}[node distance=6mm]
      \node[cfterm] (start) {start of foo};
      \node[cfstep, below=of start] (init) {let mut count = 0;};
      \node[cftest, below=of init] (test) {count < 5};
      \node[cfstep, below=8mm of test] (body) {count += 1;};
      \node[cfstep, below=of body] (after) {assert(count == 5);};
      \node[cfterm, below=of after] (end) {end of foo};

      \draw[cfflow] (start) -- (init);
      \draw[cfflow] (init) -- (test);
      \draw[cfflow] (test) -- node[cflab, right] {true} (body);
      \draw[cfflow] (body.east) -- ++(14mm, 0) |- (test.east);
      \draw[cfflow] (test.west) -- node[cflab, above] {false} ++(-14mm, 0) |- (after.west);
      \draw[cfflow] (after) -- (end);
    \end{tikzpicture}
  \end{adjustbox}
  \caption{\label{fig:flow:simple_while_loop} The flow graph of the
    \token{while} loop in Listing~\ref{lst:flow:simple_while_loop}. The
    condition is evaluated before every iteration, including the first.}
\end{figure}


\subsection{Loops entered at least once}
\label{sec:flow:loop_at_least_once}

A while loop evaluates its condition before its first iteration, so its body
may never run. Ymir has no loop that tests its condition after the body. A loop
whose body must run at least once is an infinite loop
(Section~\ref{sec:flow:inf_loop}) that breaks at the end of its body.

\begin{lstlisting}[style=coloredverbatim, caption=A loop entered at least once, label=lst:flow:loop_at_least_once]
let mut count = 0;

loop {
  count += 1;
  if count >= 5 { break; }
}

assert(count == 5);
\end{lstlisting}

\subsection{While-let loop}
\label{sec:flow:while_let_loop}

In conditional branching, as discussed in Section~\ref{sec:flow:cond_pattern}, a
pattern can serve as the condition. This capability, similar to its usage in if
conditionals, is available for while loops. In the following given example, the loop continues as long as
the option contained within the variable \token{x} holds a value and halts
either when it doesn't hold a value or when its value exceeds \token{5}. The
code presented in Listing~\ref{lst:flow:while_let_rewritten} exhibits identical
behavior, albeit implemented using a loop construct and an if let conditional,
as the compiler rewrites the code from Listing~\ref{lst:flow:while_let_example}.

\begin{lstlisting}[style=coloredverbatim, caption=While-let example, label=lst:flow:while_let_example]
fn foo(i: i32)-> i32?;

let mut x = foo(0);
let mut count = 0;

while let Ok(z) = x && z < 5 {
  count += 1;
  x = foo(count);
}
\end{lstlisting}

\begin{lstlisting}[style=coloredverbatim, caption=While-let rewritten using loop and if-let, label=lst:flow:while_let_rewritten]
fn foo(i: i32)-> i32?;

let mut x = foo(0);
let mut count = 0;

loop {
  if let Ok(z) = x && z < 5 {
    count += 1;
    x = foo(count);
  } else break;
}
\end{lstlisting}

As one might expect, since while-let is merely a syntactic transformation into a
loop over an if pattern condition, all the pattern and guard systems discussed
in Section~\ref{sec:flow:cond_pattern} are applicable to the while-let construction.

\subsection{While loop value}

While loops don't inherently produce any values, as there's no assurance that
the loop will be entered and a value constructed. Only infinite loops defined
with the \token{loop} keyword can construct a value. As a consequence, a
\textit{break} statement within a \textit{while} loop cannot be linked with a
value, and a while loop always has the type \token{void}.


\vfill%
\pagebreak
